\documentclass[10pt,a4paper]{amsart}
\usepackage[margin=19mm]{geometry}
\usepackage{amsmath,amssymb,mathtools}
\usepackage[scr=rsfs,cal=euler]{mathalfa}
\usepackage{xfrac}
\usepackage[unicode,hidelinks,bookmarks=false]{hyperref}
\newcommand{\ie}{i.e.\ }
\newcommand{\cf}{cf.\ }
\newcommand{\resp}{respectively\ }
\newcommand{\Subsection}{\subsection}
\providecommand{\nobreakdash}{\nobreak\hbox{-}}
\setlength{\emergencystretch}{2em}
\multlinegap0pt
\usepackage{xcolor}
\usepackage{amssymb}
\usepackage{enumerate}
\usepackage{float}
\usepackage{dsfont}
\newtheorem{theorem}{Theorem}
\newtheorem{lemma}{Lemma}
\DeclareMathOperator{\C}{\mathbb{C}}
\DeclareMathOperator{\D}{\mathbb{D}}
\title{Eigenvalues for Infinitesimal Generators of Semigroups of Composition Operators}
\author{Maria Kourou, Eleftherios K. Theodosiadis, Konstantinos Zarvalis}
\date{Source: arXiv:2508.00666v2 (CC BY 4.0)}
\begin{document}
\maketitle

\noindent\emph{Source and licence.} Eigenvalues for Infinitesimal Generators of Semigroups of Composition Operators. Version arXiv:2508.00666v2 (CC BY 4.0), distributed by arXiv under the Creative Commons Attribution 4.0 International licence, http://creativecommons.org/licenses/by/4.0/, as recorded in the arXiv licence metadata for this exact version and shown on the abstract page https://arxiv.org/abs/2508.00666v2. Full licence notice: \texttt{licenses/arxiv-2508.00666-CC-BY-4.0.txt}.\par\smallskip

\noindent\textit{Editorial scope.} Counted unit: the article's theorem dirichlet (source0.tex lines 189--200). The article's complete original proof of it is delivered with it (source0.tex lines 624--654, one \texttt{proof} environment of 31 lines, which the article places away from the statement). No mathematics is written, repaired or re-derived: the statement and the proof are the article's own lines, in the article's own notation, and no retained line is substituted or re-flowed.  Retained uncounted as the source-local context the proof needs: lines 169--181; lines 530--532; lines 576--581.  Nothing was omitted from any retained span, so the package carries no omission record.  No retained line contains a citation, so the wrapper carries no bibliography and no bibliography entry.  No retained line contains a figure environment, an image-inclusion command or drawing code, so no image is delivered and the package declares no asset dependency.  Editorial formatting only, in addition: an AMS \texttt{amsart} preamble that restates the article's own declarations and macros used by the retained lines, namely the environments \texttt{theorem}, \texttt{lemma} on separate counters, together with the standard packages those lines need.  The retained environments print in this document as Theorem 1, Lemma 1 in order of appearance, whereas the article prints the same items with the numbering of its own full document.  The complete native source of the article as distributed in the arXiv e-print for this exact version is bundled unchanged as \texttt{sources/182560/source0.tex}.
	\begin{theorem}\label{thm:intro defining function}
		Let $(\phi_t)$ be a hyperbolic semigroup in $\D$ which induces the semigroup of composition operators $(T_t)$. Suppose that $\Omega$ is the Koenigs domain of $(\phi_t)$ and $\psi_\Omega:(a,b)\to[-\infty,+\infty)$ its defining function, where $-\infty<a<b<+\infty$. If $\sigma_\mathcal{D}$ is the point spectrum of the infinitesimal generator of $(T_t)$ when acting on the Dirichlet space, then:
		\begin{enumerate}
			\item[\textup{(a)}] It holds that $\lambda\in\sigma_\mathcal{D}\setminus\{0\}$ if and only if $\int_{a}^{b}e^{2\mathrm{Re}\lambda\psi_\Omega(y)}dy<+\infty$.
		\end{enumerate}
		Set $c_\Omega=\inf\{c\in(-\infty,0]:\int_{a}^{b}e^{2c\psi_\Omega(y)}dy<+\infty\}$. Then:
		\begin{enumerate}
			\item[\textup{(b)}]  If $c_\Omega=0$, then $\sigma_\mathcal{D}=\{0\}$.
			\item[\textup{(c)}] If $c_\Omega\in(-\infty,0)$ and $\int_{a}^{b}e^{2c_\Omega\psi_\Omega(y)}dy=+\infty$, then $\sigma_\mathcal{D}=\{\lambda\in\C:c_\Omega<\mathrm{Re}\lambda<0\}\cup\{0\}$.
			\item[\textup{(d)}] If $c_\Omega\in(-\infty,0)$ and $\int_{a}^{b}e^{2c_\Omega\psi_\Omega(y)}dy<+\infty$, then $\sigma_\mathcal{D}=\{\lambda\in\C:c_\Omega\le\mathrm{Re}\lambda<0\}\cup\{0\}$.
			\item[\textup{(e)}] If $c_\Omega=-\infty$, then $\sigma_\mathcal{D}=\{\lambda\in\C:-\infty<\mathrm{Re}\lambda<0\}\cup\{0\}$.
		\end{enumerate}
	\end{theorem}

	\begin{lemma}\label{lm:hyperbolic basic condition}
		Let $(\phi_t)$ be a hyperbolic semigroup in $\D$ which induces the semigroup of composition operators $(T_t)$. Suppose that $\Omega$ is the Koenigs domain of $(\phi_t)$ and that $\sigma_\mathcal{D}$ is the point spectrum of the infinitesimal generator of $(T_t)$. Then $\{0\}\subseteq\sigma_{\mathcal{D}}\subseteq iS(0,\pi)\cup\{0\}$. In particular, both inclusions are sharp.
	\end{lemma}

	\begin{lemma}\label{lm:hyperbolic x parametrization}
		Let $(\phi_t)$ be a hyperbolic semigroup in $\D$ which induces the semigroup of composition operators $(T_t)$. Suppose that $\Omega$ is the Koenigs domain of $(\phi_t)$ and that $\sigma_\mathcal{D}$ is the point spectrum of the infinitesimal generator of $(T_t)$. Then $\lambda\in\sigma_{\mathcal{D}}\setminus\{0\}$ if and only if
		\begin{equation}
			\int\limits_{-\infty}^{0}e^{2\mathrm{Re}\lambda x}\ell_\Omega(x)dx<+\infty.
		\end{equation}
	\end{lemma}

	\begin{theorem}\label{thm:intro hyperbolic dirichlet}
		Let $(\phi_t)$ be a hyperbolic semigroup in $\D$ which induces the semigroup of composition operators $(T_t)$. Suppose that $\Omega$ is the Koenigs domain of $(\phi_t)$ and $\sigma_\mathcal{D}$ the point spectrum of the infinitesimal generator of $(T_t)$. Then:
		\begin{enumerate}
			\item[\textup{(a)}] If $W(\Omega)=+\infty$, then $\sigma_{\mathcal{D}}=\{0\}$.
			\item[\textup{(b)}] If $W(\Omega)<+\infty$ and $\delta_\Omega=+\infty$, then $\sigma_{\mathcal{D}}=\{\lambda\in\C:\mathrm{Re}\lambda<0\}\cup\{0\}$.
			\item[\textup{(c)}] If $W(\Omega)<+\infty$, $\delta_\Omega\in(0,+\infty)$, and $\int_{-\infty}^{0}e^{-2\delta_\Omega x}\ell_\Omega(x)dx<+\infty$, then
			$$\sigma_{\mathcal{D}}=\{\lambda\in\C:-\delta_\Omega\le\mathrm{Re}\lambda<0\}\cup\{0\}.$$
			\item[\textup{(d)}] If $W(\Omega)<+\infty$, $\delta_\Omega\in(0,+\infty)$, and $\int_{-\infty}^{0}e^{-2\delta_\Omega x}\ell_\Omega(x)dx=+\infty$, then
			$$\sigma_{\mathcal{D}}=\{\lambda\in\C:-\delta_\Omega<\mathrm{Re}\lambda<0\}\cup\{0\}.$$
			\item [\textup{(e)}] If $W(\Omega)<+\infty$ and $\delta_\Omega=0$, then $\sigma_{\mathcal{D}}=\{0\}$.
		\end{enumerate}
	\end{theorem}

	\begin{proof}[Proof of Theorem \ref{thm:intro hyperbolic dirichlet}]
		(a) Certainly $0\in\sigma_{\mathcal{D}}$. Suppose that $\lambda\in\sigma_{\mathcal{D}}\setminus\{0\}$. By Lemma \ref{lm:hyperbolic basic condition}, we have $\mathrm{Re}\lambda<0$ and hence $e^{2\mathrm{Re}\lambda x}>1$, for all $x\in(-\infty,0)$. Therefore
		$$\int\limits_{-\infty}^{0}e^{2\mathrm{Re}\lambda x}\ell_\Omega(x)dx\ge\int\limits_{-\infty}^{0}\ell_\Omega(x)=W(\Omega)=+\infty.$$
		However, due to Lemma \ref{lm:hyperbolic x parametrization}, we obtain $\lambda\notin\sigma_{\mathcal{D}}$. Contradiction! This leads to $\sigma_{\mathcal{D}}\setminus\{0\}=\emptyset$ and the desired outcome.
		
		(b) Because of Lemma \ref{lm:hyperbolic basic condition}, it suffices to prove that $iS(0,\pi)\subseteq\sigma_{\mathcal{D}}$. Towards this goal, let $\lambda\in\C$ with $\mathrm{Re}\lambda<0$. Pick $M>0$ such that $-M<\mathrm{Re}\lambda$ or equivalently $M+\mathrm{Re}\lambda>0$. Due to our hypothesis, for this $M$ we may find some $x_0\in(-\infty,0)$ such that $\log\ell_\Omega(x)/(2x)>M$, for all $x\in(-\infty,x_0)$. Elementary rearrangements show that $\ell_\Omega(x)<e^{2Mx}$, for all $x\in(-\infty,x_0)$. Then
		\begin{eqnarray*}
			\int\limits_{-\infty}^{0}e^{2\mathrm{Re}\lambda x}\ell_\Omega(x)dx &=& \int\limits_{-\infty}^{x_0}e^{2\mathrm{Re}\lambda x}\ell_\Omega(x)dx+\int\limits_{x_0}^{0}e^{2\mathrm{Re}\lambda x}\ell_\Omega(x)dx \\
			&\le&\int\limits_{-\infty}^{x_0}e^{2(\mathrm{Re}\lambda+M)x}dx+\int\limits_{x_0}^{0}e^{2\mathrm{Re}\lambda x}\ell_\Omega(x)dx\\
			&<&+\infty,
		\end{eqnarray*}
		since on one hand, $2(\mathrm{Re}\lambda+M)>0$ and on the other hand, the integral from $x_0$ to $0$ is clearly finite. By Lemma \ref{lm:hyperbolic x parametrization} we conclude that $\lambda\in\sigma_{\mathcal{D}}$. The arbitrariness in the choice of $\lambda\in iS(0,\pi)$ implies the desired inclusion.
		
		(c) Due to our assumptions and combining Lemma \ref{lm:hyperbolic x parametrization} with the vertical invariance of $\sigma_\mathcal{D}$, we deduce that $\{w\in\C:\mathrm{Re}w=-\delta_\Omega\}\subseteq\sigma_{\mathcal{D}}$. The shape of the point spectrum imposed by Theorem \ref{thm:intro defining function} shows that $\widetilde{\Sigma}_{\delta_\Omega}\subseteq\sigma_{\mathcal{D}}$. Clearly $0\in\sigma_{\mathcal{D}}$ as well. For the reverse inclusion, assume that $\lambda\in\C$ with $\mathrm{Re}\lambda<-\delta_\Omega$. We are going to prove that $\lambda\notin\sigma_{\mathcal{D}}$. Fix $\epsilon\in(0,-\mathrm{Re}\lambda-\delta_\Omega)$. Due to the assumption on the limit, for this $\epsilon$ we may find $x_1\in(-\infty,0)$ such that $\log\ell_\Omega(x)/(2x)<\delta_\Omega+\epsilon$, for all $x\in(-\infty,x_1)$. Quick calculations yield $\ell_\Omega(x)>e^{2(\delta_\Omega+\epsilon)x}$, for all $x\in(-\infty,x_1)$. Therefore,
		\begin{equation*}
			\int\limits_{-\infty}^{0}e^{2\mathrm{Re}\lambda x}\ell_\Omega(x)dx\ge\int\limits_{-\infty}^{x_1}e^{2(\mathrm{Re}\lambda+\delta_\Omega+\epsilon)x}dx=+\infty,
		\end{equation*}
		since by construction $2(\mathrm{Re}\lambda+\delta_\Omega+\epsilon)<0$. Returning to Lemma \ref{lm:hyperbolic x parametrization} we understand that $\lambda\notin\sigma_{\mathcal{D}}$. Combining with the fact that $\sigma_{\mathcal{D}}\subseteq iS(0,\pi)\cup\{0\}$, we get the required equality for the point spectrum.
		
		(d) Working identically as in the previous case, we may reach the point where $\sigma_{\mathcal{D}}\subseteq \widetilde{\Sigma}_{\delta_\Omega}\cup\{0\}$. Due to the integral condition in our hypothesis, we further get that $\sigma_{\mathcal{D}}\subseteq \Sigma_{\delta_\Omega}\cup\{0\}$. Reversely, let $\lambda\in\C$ with $\mathrm{Re}\lambda\in(-\delta_\Omega,0)$. Fix $\epsilon\in(0,\mathrm{Re}\lambda+\delta_\Omega)$. For this $\epsilon$ we may find $x_2\in(-\infty,0)$ such that $\log\ell_\Omega(x)/(2x)>\delta_\Omega-\epsilon$, for all $x\in(-\infty,x_2)$, due to the existence of the limit. Again, we can rearrange to obtain $\ell_\Omega(x)<e^{2(\delta_\Omega-\epsilon)x}$, for all $x\in(-\infty,x_2)$. Then
		\begin{equation*}
			\int\limits_{-\infty}^{0}e^{2\mathrm{Re}\lambda x}\ell_\Omega(x)dx\le\int\limits_{x_2}^{0}e^{2\mathrm{Re}\lambda x}\ell_\Omega(x)dx+\int\limits_{-\infty}^{x_2}e^{2(\mathrm{Re}\lambda+\delta_\Omega-\epsilon)x}dx<+\infty,
		\end{equation*}
		since $2(\mathrm{Re}\lambda+\delta_\Omega-\epsilon)>0$, while the integral from $x_2$ to $0$ is evidently finite. Returning to Lemma \ref{lm:hyperbolic x parametrization}, we see that $\lambda\in\sigma_{\mathcal{D}}$. Combining with the fact that $0$ is always an eigenvalue, we deduce that $\Sigma_{\delta_\Omega}\cup\{0\}\subseteq\sigma_{\mathcal{D}}$ which produces the equality.
		
		(e) By Lemma \ref{lm:hyperbolic basic condition} it suffices to show that any $\lambda\in iS(0,\pi)$ is not an eigenvalue. Fix such a $\lambda$ with $\mathrm{Re}\lambda<0$ and find $\epsilon\in(0,-\mathrm{Re}\lambda)$. For this $\epsilon$, our assumption on the limit implies the existence of some $x_3\in(-\infty,0)$ so that $\log\ell_\Omega(x)/(2x)<\epsilon$, for all $x\in(-\infty,x_3)$. Similarly to before, this means that $\ell_\Omega(x)>e^{2\epsilon x}$, for all $x\in(-\infty,x_3)$. Consequently,
		\begin{equation*}
			\int\limits_{-\infty}^{0}e^{2\mathrm{Re}\lambda x}\ell_\Omega(x)dx\ge\int\limits_{-\infty}^{x_3}e^{2(\mathrm{Re}\lambda+\epsilon)x}dx=+\infty,
		\end{equation*}
		because $2(\mathrm{Re}\lambda+\epsilon)<0$. Once again, Lemma \ref{lm:hyperbolic x parametrization} yields the desired outcome.
	\end{proof}

\end{document}
